\section{De MoE a Attention: por qué la innovación en arquitectura vuelve a ser importante}

Esta sección sitúa a Attention dentro de una historia más amplia de las arquitecturas. El Transformer contiene Attention, FFN y codificación posicional. La FFN ya fue transformada por MoE; la codificación posicional pasó de posiciones absolutas a posiciones relativas, RoPE, etc.; Attention, por su parte, sigue evolucionando por vías como Full, Linear, Sparse y Sliding Window. La tesis central de la entrevista es que MoE ha sido el mayor avance de arquitectura de los últimos años y que la siguiente pieza podría ser Attention.

\begin{figure}[H]
\centering
\includegraphics[width=0.96\textwidth]{moe-to-attention.png}
\caption{De MoE a Attention: la FFN ya fue esculpida en forma de MoE, y Attention se convierte en la siguiente región por esculpir. Diagrama conceptual de elaboración propia, a partir de la conversación de 00:58:08--01:02:48.}
\end{figure}

\begin{knowledgebox}{Lectura de la figura: el Transformer no es un bloque monolítico}
El Transformer es un apilamiento de Attention + FFN. Al convertirse la FFN en MoE, cambió la forma de usar el cómputo; Attention sigue siendo el cuello de botella del contexto largo, por lo que vías como Linear/Sparse intentan seguir «esculpiendo» la arquitectura.
\end{knowledgebox}

\subsection{FLOPs, loss y arquitectura}

Esta subsección convierte la elección de arquitectura en una función objetivo comparable. Las empresas de modelos no modifican la arquitectura por novedad, sino para obtener mejores resultados de entrenamiento e inferencia con el mismo presupuesto de cómputo.

El objetivo elemental detrás de la innovación en arquitectura es: con los mismos FLOPs, obtener una loss más baja. FLOPs significa floating point operations, es decir, la cantidad de operaciones de punto flotante. Una buena arquitectura no se limita a reducir el cómputo, sino que, con el mismo presupuesto, expresa funciones con mayor eficacia, aprovecha mejor los datos y se adapta con más eficiencia al hardware.

\begin{importantbox}{La función objetivo de la innovación en arquitectura}
Bajo las mismas restricciones de datos y cómputo, la innovación en arquitectura reduce el costo de entrenamiento/inferencia, aumenta la capacidad expresiva o mejora la eficiencia en hardware; lo ideal es mejorar las tres a la vez.
\end{importantbox}

\subsection{Por qué el agotamiento de los datos hace más importante la innovación algorítmica}

A continuación, se devuelven los algoritmos al contexto de las restricciones de la industria. Mientras los datos y el cómputo puedan seguir creciendo sin mayor reflexión, la importancia de la innovación en arquitectura queda fácilmente oculta; cuando ambos se encarecen, la eficiencia algorítmica vuelve al centro.

Al inicio del programa se menciona que los datos, los algoritmos y el cómputo son los tres motores que impulsan la inteligencia artificial. Cuando el crecimiento de los datos se desacelera y el cómputo se vuelve limitado o costoso, el valor de los algoritmos y de la arquitectura vuelve a aumentar. Las empresas chinas, con un cómputo más limitado que el de Estados Unidos, tienen precisamente por ello más incentivos para innovar en direcciones como MoE, Sparse Attention, Linear Attention y NoPE/RoPE, buscando «ahorrar cómputo sin perder rendimiento».

\begin{knowledgebox}{Cambio en las restricciones de los tres motores}
\begin{tabularx}{\textwidth}{p{0.18\textwidth}p{0.38\textwidth}X}
\toprule
Variable & Práctica principal en el pasado & Presión actual \\
\midrule
Datos & Seguir ampliando los datos de internet/código/multimodales & El crecimiento de los datos públicos de alta calidad se desacelera. \\
Cómputo & Acumular GPU y ampliar la escala de entrenamiento & Se agravan las restricciones de costo, suministro y energía. \\
Algoritmos & Mantener la base del Transformer con ajustes menores & Se necesita mayor eficacia con los mismos FLOPs. \\
\bottomrule
\end{tabularx}
\end{knowledgebox}

\subsection{Resumen de la sección}

Attention vuelve a ser el foco porque el contexto largo y el costo de inferencia llevan a la atención global a una posición de cuello de botella. MoE transformó la FFN; la siguiente etapa podría ser que Linear/Sparse/Hybrid Attention transforme la capa de atención.
